\phantomsection
\chapter*{Kavramlar Sözlüğü}
\markboth{Kavramlar Sözlüğü}{Kavramlar Sözlüğü}
\addcontentsline{toc}{chapter}{Kavramlar Sözlüğü}

\begin{multicols}{2}
\raggedcolumns
\small

\textbf{Anlamlılık düzeyi ($\alpha$).} Testten önce belirlenen Tip I hata
eşiği.\par\smallskip

\textbf{Bağımsızlık.} Bir gözlemin veya değişkenin diğerine ilişkin ek bilgi
taşımaması durumu; anlamı kullanılan modele göre belirlenir.\par\smallskip

\ifimoders\else
\textbf{Bootstrap.} Gözlenen örneklemden yerine koyarak yeniden örnekler
seçip bir tahminin örnekleme değişkenliğini yaklaşık olarak inceleyen yöntem.\par\smallskip
\fi

\textbf{Çapraz tablo.} İki kategorik değişkenin ortak frekanslarını gösteren
tablo.\par\smallskip

\textbf{Etki büyüklüğü.} Farkın veya ilişkinin uygulamalı büyüklüğünü ortak
bir ölçekte özetleyen nicelik.\par\smallskip

\textbf{Eksik veri.} Araştırmada tanımlanan bir değişkenin belirli bir analiz
birimi için gözlenmemiş olması; sıfır veya uygulanamaz değerle aynı değildir.\par\smallskip

\textbf{Evren.} Araştırma sonucunun genellenmek istendiği birimlerin tümü.\par\smallskip

\textbf{Güven aralığı.} Tekrarlı örneklemede belirli bir kapsama oranına sahip
olacak yöntemle üretilen parametre aralığı.\par\smallskip

\textbf{İstatistik.} Örneklem verilerinden hesaplanan sayısal özet.\par\smallskip

\textbf{Merkezi Limit Teoremi.} Uygun koşullarda, örneklem hacmi büyüdükçe
standartlaştırılmış örneklem ortalamasının dağılımının normale yaklaşması.\par\smallskip

\textbf{MCAR.} Eksik kalma olasılığının gözlenen ve gözlenmeyen değerlerle
ilişkili olmadığı tamamen rassal eksiklik varsayımı.\par\smallskip

\textbf{MAR.} Eksikliğin, analize alınan gözlenen bilgiler koşullu olarak
gözlenmeyen değere ayrıca bağlı olmadığı varsayım.\par\smallskip

\textbf{MNAR.} Gözlenen bilgiler hesaba katıldıktan sonra bile eksikliğin
gözlenmeyen değer veya kaydedilmemiş nedenle ilişkili olması.\par\smallskip

\textbf{Nedensellik.} Bir değişkendeki müdahalenin başka bir değişkende
değişim üretmesi; yalnız ilişki ölçüleriyle kurulamaz.\par\smallskip

\textbf{Nokta tahmini.} Bilinmeyen parametre için örneklemden elde edilen tek
sayısal değer.\par\smallskip

\textbf{Rassal değişken.} Olasılıksal bir sürecin sonuçlarına sayısal değer
atayan fonksiyon.\par\smallskip

\textbf{Örneklem.} Evren hakkında bilgi edinmek üzere gözlenen birimler
kümesi.\par\smallskip

\textbf{Örnekleme dağılımı.} Bir istatistiğin tekrarlı örneklemlerde alacağı
değerlerin olasılık dağılımı.\par\smallskip

\textbf{Örnekleme değişkenliği.} Aynı süreçle seçilen farklı örneklemlerin
farklı istatistikler üretmesi.\par\smallskip

\textbf{\pval-değeri.} $H_0$ doğru kabul edildiğinde gözlenen kadar veya daha uç
bir test sonucu elde etme olasılığı.\par\smallskip

\textbf{Parametre.} Evrene veya istatistiksel modele ait sabit fakat genellikle
bilinmeyen nicelik.\par\smallskip

\textbf{Regresyon eğimi.} Açıklayıcı değişkendeki bir birimlik artışla
ilişkili beklenen yanıt değişimi.\par\smallskip

\ifimoders\else
\textbf{Rastgeleleştirme testi.} Sıfır hipotezi altında değiştirilebilir olan
etiketleri veya işaretleri yeniden düzenleyerek test dağılımı oluşturan yöntem.\par\smallskip
\fi

\textbf{Serbestlik derecesi.} Bir istatistiğin dağılımını belirleyen bağımsız
bilgi miktarı.\par\smallskip

\textbf{Standart hata.} Bir tahmin edicinin örnekleme dağılımındaki standart
sapma.\par\smallskip

\textbf{Standart sapma.} Gözlemlerin dağılım içindeki değişkenliğini, kendi
biriminde özetleyen yayılım ölçüsü. Standart hatayla aynı nicelik değildir.\par\smallskip

\textbf{Tahmin.} Bir tahmin edicinin gözlenen örneklemde aldığı sayısal değer.\par\smallskip

\textbf{Tahmin edici.} Örneklemden parametre tahmini üreten, örneklem
seçilmeden önce rassal olan kural.\par\smallskip

\textbf{Tip I hata.} Doğru bir $H_0$ hipotezini reddetme.\par\smallskip

\textbf{Tip II hata.} Yanlış bir $H_0$ hipotezini reddedememe.\par\smallskip

\textbf{Çoklu atama.} Eksik değerler için birden çok tamamlanmış veri seti
üretip analiz tahminleri ile atama belirsizliğini birlikte birleştiren yöntem.\par\smallskip

\textbf{Test gücü.} Belirli bir alternatif doğruyken $H_0$'ı reddetme
olasılığı; $1-\beta$ ile gösterilir.\par\smallskip

\textbf{Yanlılık.} Tahminin veya veri üretim sürecinin hedef değerden
sistematik olarak sapması.\par\smallskip

\ifimoders\else
\textbf{Varyans büyütme faktörü (VIF).} Bir regresyon katsayısının
belirsizliğinin açıklayıcılar arası doğrusal ilişkiler nedeniyle ne ölçüde
büyüdüğünü özetleyen tanı ölçüsü.\par\smallskip

\textbf{Varyans analizi (ANOVA).} Bir nicel yanıtın bağımsız gruplardaki
ortalamalarını, gruplar arası değişkenliği grup içi değişkenliğe oranlayarak
karşılaştıran yöntem ailesi.\par\smallskip

\textbf{Welch ANOVA.} Grup varyanslarının eşitliğini gerektirmeyen, yaklaşık
payda serbestlik derecesi kullanan genel ortalama testi.\par\smallskip

\textbf{Eta-kare ($\eta^2$).} ANOVA'da toplam örneklem değişkenliğinin
gruplar arası farklılıkla ilişkili oranı.\par\smallskip

\textbf{Omega-kare ($\omega^2$).} Eta-karenin küçük örneklemdeki yukarı
yönlü yanlılığını azaltmayı amaçlayan ANOVA etki büyüklüğü.\par\smallskip
\fi

\textbf{Yansız tahmin edici.} Beklenen değeri hedef parametreye eşit olan
tahmin edici.\par\smallskip

\textbf{$R^2$.} Regresyon modelinin örneklemde açıkladığı yanıt değişkenliği
oranı.\par\smallskip

\normalsize
\end{multicols}